\section{Dados e recorte da análise}
\label{sec:dados}

\subsection{Fontes, unidades e preparação}

A base econômica é a MIP brasileira de 2015, disponibilizada pelo IBGE no detalhamento de 67 atividades \parencite{ibge2015mip}. A atividade produtiva é a unidade de análise, e os valores monetários são expressos em milhões de reais, a preços básicos de 2015.

\textcite{sanguinetAzzoni2024} fornecem intensidades diretas de CO\textsubscript{2} por unidade de produção bruta para 2011 e 2018. Para o cenário de 2015, utiliza-se interpolação linear. A interpolação pressupõe evolução linear da intensidade entre os dois anos e não incorpora observação física de emissões em 2015.

\subsection{Fronteira e validade ambiental}

O recorte abrange emissões de CO\textsubscript{2} das atividades produtivas brasileiras para demanda doméstica e exportações. Exclui emissões estrangeiras na produção de importações e no processamento ou uso posterior dos produtos exportados. O destino da atribuição identifica o setor do produto final na fronteira nacional. A conta pela demanda não corresponde, portanto, à pegada mundial do consumo dos residentes.

A análise contabiliza CO\textsubscript{2}, sem agregar os demais gases de efeito estufa em CO\textsubscript{2} equivalente. A cobertura ambiental permanece incerta: a documentação da fonte apresenta descrições divergentes sobre uso da terra, e não foram recuperados os numeradores físicos nem os procedimentos completos de alocação das emissões. A correspondência setorial não certifica a cobertura nem o volume físico atribuído a cada atividade.

A compatibilidade monetária também permanece pendente. A fonte relata atualização de sua matriz inter-regional para preços de dezembro de 2018, enquanto esta aplicação combina coeficientes interpolados com valores de 2015, sem ajuste adicional. A igualdade das unidades não resolve essa diferença de preços; é necessário verificar os denominadores setoriais antes de interpretar os volumes estimados.

A base ambiental constitui, assim, um \textbf{cenário exploratório}. As verificações de balanço e das identidades contábeis asseguram consistência dos cálculos, mas não validam as emissões físicas. Conclusões ambientais definitivas dependem da reconstrução e validação dos coeficientes ou de sua substituição por uma extensão compatível com 2015.
